\scanpage{10}
\textcircled{2.5} Situation où $N$ invariant \unsure{sous|dans} $\unsure{\mathcal E|\mathscr E}$ i.e. inv. par $\unsure{\mathcal G|\mathscr G}$, et $\mathcal G\mathbin{-}\Gamma$ est surjectif : $\deleted{\unsure{\mathcal E}}\mathcal G'$ est \deleted{\unsure{décrit|alors|devient}} muni de \unsure{suite|suites} de composition ayant la structure suivante (\unsure{quots|quotients} successifs $\Gamma,\unsure{\mathfrak G|G|\mathscr G},G$)
\diagram{(24)\quad\begin{tikzcd}[ampersand replacement=\&,column sep=small,row sep=tiny]
{}\&{}\&\deleted{\unsure{\mathcal E}}\mathcal G'\&{}\&{}\\
G\arrow[rrrr,dash,start anchor=north,end anchor=north,to path={-- ++(0,.17) -| (\tikztotarget)}]\arrow[rr,dash,start anchor=south,end anchor=south,to path={-- ++(0,-.18) -| (\tikztotarget)}]\&{}\&\mathfrak G\arrow[rr,dash,start anchor=south,end anchor=south,to path={-- ++(0,-.18) -| (\tikztotarget)}]\&{}\&\Gamma\\
{}\&\mathfrak G'\&{}\&\mathcal G\&{}
\end{tikzcd}\quad.}

L’extension $\mathfrak G'$ de $\mathfrak G$ par $G=\mathrm{Aut}(\unsure{\mathscr X|\mathcal X|\mathfrak X}'/\mathscr X)$ est de ``nature géométrique'', l’extension de $\Gamma$ par $\mathfrak G'$ exprime les caractères arithmétiques de la situation.

Nous nous intéresserons surtout à la situation où $\mathfrak G$ et $G$ sont finis, i.e. $\mathfrak G'$ fini, et $\Gamma$ infini, i.e. $\bar k$ extension infinie de $k$, \deleted{Considérons} $\mathscr X$ \deleted{\unsure{étant}} (donc aussi $\mathscr X'$) étant de prés. finie sur $k$. [On a envie de plus, $\mathscr X/\bar k/k$ étant fixés, de faire varier $\mathscr X'$ \inserted{ou $N$} de façon que $\widetilde{\mathscr X}$ soit la lim proj. des $\mathscr X'$ i.e. $\bigcap N=1$.] On voit aisément qu’on peut trouver alors un groupe quotient fini $\deleted{\unsure{k}}\gamma_i$ \inserted{$(=\Gamma/\Gamma_i)$} de $\Gamma$
% archive sheet number: 9
